\chapter{嵌入式安全生命周期}
\label{chap:security-lifecycle}

嵌入式安全不是在启动链上增加一次签名校验，而是从芯片信任根、生产制造、系统启动、运行隔离、远程更新到退役的持续工程活动。
安全设计必须建立在明确的资产、攻击者能力和信任边界上，并与可恢复性、可维护性和密钥生命周期共同设计。

产品安全还具有时间维度：发布时未被利用的缺陷，可能在供应链披露、新攻击技术出现或云端权限变化后变成可利用漏洞。工程团队因此需要同时交付安全需求、实现、测试证据、组件清单、更新能力、漏洞响应流程和支持期限。NIST SSDF 把安全活动嵌入软件开发过程\cite{nist-ssdf}；面向联网设备的能力基线还要求设备识别、配置、数据保护、接口控制、更新和状态感知等可维护能力\cite{nist-iot-core-baseline}。

本章关注产品生命周期和组织闭环。硬件根身份、TEE、RPMB 与远程证明的机制见 \mbox{第~\ref{chap:trusted-execution-identity}~章}，构建与 A/B 更新见 \mbox{第~\ref{chap:system-build-update}~章}，持久化和真实掉电恢复见 \mbox{第~\ref{chap:storage-reliability}~章}。

\section{威胁模型与安全目标}

威胁模型至少应回答：保护什么、攻击者能接触什么、哪些组件被信任、攻击成功会造成什么影响。
常见资产包括固件完整性、设备身份、用户数据、控制权限、算法和云端凭据。

攻击者能力必须写成工程假设，而不是统一称为“黑客”。常见层次包括：仅能从互联网发包的远程攻击者、处于同一局域网或无线覆盖范围的邻近攻击者、获得普通用户权限的本地软件、能够拆机连接调试口的物理攻击者，以及能够影响源码、依赖、构建或生产工位的供应链攻击者。不同能力对应不同成本和防护边界。

嵌入式设备面临的攻击面包括：
\begin{itemize}
  \item 网络服务、无线接口和管理协议；
  \item 调试口、启动接口、可移动存储和测试点；
  \item 更新服务器、构建系统、生产线和供应链；
  \item 本地文件系统、IPC、驱动 ioctl 和解析器；
  \item 电压、时钟、温度和物理篡改形成的故障注入。
\end{itemize}
安全目标应可验证，例如“设备只执行由产品发布密钥授权且版本不低于安全计数器的固件”，而不是笼统地写成“系统安全”。

\begin{table}[H]
  \centering
  \caption{威胁模型需要固定的工程输入}
  \begin{tabular}{p{27mm}p{40mm}p{42mm}p{27mm}}
    \hline
    输入 & 需要说明 & 容易遗漏的例子 & 产出 \\
    \hline
    资产 & 机密性、完整性、可用性和新鲜性 & 校准值、启动计数、云端撤权状态 & 资产清单与责任人 \\
    攻击者能力 & 访问距离、权限、设备时间和预算 & 邻近无线、拆机、生产工位账户 & 能力与排除假设 \\
    信任边界 & 谁验证谁、密钥和数据跨越何处 & 处理器域、安全世界与普通世界、设备与云端 & 数据流与边界图 \\
    失败影响 & 安全、隐私、业务和现场恢复代价 & 批量失联、错误控制、不可返修 & 风险等级与处置门槛 \\
    \hline
  \end{tabular}
\end{table}

威胁模型应绑定产品型号、硬件修订、软件基线和部署形态。加入蓝牙配网、开放 USB 维护口、更换云端身份提供方或把设备部署到公共网络，都可能改变原有攻击面，必须触发增量评审。被排除的威胁也要记录理由和有效条件，例如“攻击者无法长期物理接触”不适用于无人值守室外设备。

\section{安全需求、风险与发布门槛}

每个高价值资产都应映射到预防、检测、响应和恢复要求。签名验证属于预防；启动失败计数和审计属于检测；撤销密钥和停止发布属于响应；可信恢复镜像属于恢复。只有预防而没有恢复，可能把单次异常变成永久变砖；只有检测而没有处置责任人，则告警不会降低风险。

安全需求应包含明确边界和判定器，例如：

\begin{itemize}
  \item 未认证网络端点不能调用改变物理输出的命令；
  \item 维护授权只能绑定单台设备、规定操作和有限时间窗口；
  \item 发布镜像必须可追溯到经评审源码、锁定依赖和指定构建环境；
  \item 已撤销证书不能在离线宽限期之外继续获得云端控制权限；
  \item 安全更新失败后必须保留一个不低于安全下限的可启动版本。
\end{itemize}

CVSS 用于描述漏洞技术严重度，但不能独立决定产品风险\cite{first-cvss-v4}。实际决策还要考虑功能是否编译启用、攻击路径是否可达、所需权限、设备部署规模、可检测性、安全或隐私后果、补丁副作用和现场更新成功率。团队应记录接受、缓解、修复或停止支持的决定、证据和批准人。

发布门槛应区分阻断项与可跟踪债务。签名链失效、默认共享凭据、可远程未认证执行、密钥进入镜像、已知被利用漏洞和无法恢复的更新缺陷通常应阻断发布；低风险诊断信息暴露可以在有责任人、截止期和补偿控制时进入债务清单。

\section{硬件信任根}

信任根是启动链最初不可被普通软件替换的验证基础，通常由 Boot ROM、一次性可编程存储、硬件唯一密钥、安全元件或受保护执行环境构成。

信任根必须小、稳定且具有可恢复的上层策略。把网络协议栈、复杂文件系统或产品业务放入不可更新区域，会扩大永久漏洞风险。Boot ROM 应只承担必要的镜像定位、基本硬件初始化和下一阶段认证，并允许通过受保护的可更新密钥层扩展信任，而不是永久绑定单一算法和单一发布密钥。

\subsection{设备身份与密钥派生}

设备唯一密钥不应以明文形式出现在通用镜像、日志或生产脚本中。
优先在受保护硬件内生成或派生设备密钥，并通过证书或注册流程建立设备身份。

不同用途的密钥应通过独立派生上下文隔离，例如启动验证、存储加密、设备认证和调试授权不能直接复用同一密钥。

设备身份还需要明确“谁为它背书”。芯片唯一值只能证明某颗芯片，制造证书可以绑定板卡和产品序列号，云端注册再绑定租户、所有权和服务权限。MAC 地址、可读序列号或应用配置都可以被复制，不能单独作为不可克隆身份。

\subsection{熔丝与生命周期状态}

安全熔丝、调试锁和防回滚计数器通常不可逆。烧写前必须验证芯片型号、板卡身份、目标值、电源条件和当前生命周期状态，并保留审计记录。
开发、生产、返修和退役阶段应具有不同权限；量产设备不能长期保留无认证的调试入口。

生命周期状态转换应由设备自身验证前置条件：生产锁定前确认恢复镜像、设备证书、随机数质量和量产测试全部有效；维修解锁需要认证授权；退役状态拒绝重新注册为正常设备。工位只发出请求而不能单方面宣称状态成功，最终结果应从目标芯片回读并写入追溯记录。

\subsection{调试授权}

完全关闭 JTAG/SWD 可以缩小攻击面，却也可能失去故障分析能力。需要保留调试的产品应使用 challenge-response 或签名授权，使授权绑定设备身份、允许的调试能力、时间或启动次数和工单。量产主密钥或通用“工程密码”不能散落在返修工具中。

调试授权失败必须保持默认锁定，且不能因实时时钟未同步而无限开放。离线返修可以使用单调计数、一次性 nonce 或短期离线授权包，但应防止旧授权在另一台设备或另一阶段重放。

\section{可信启动链}

可信启动由前一级验证后一级，直到内核、设备树、根文件系统或关键应用。
每一级都应验证载荷身份、完整性、版本和适用设备，并在失败时进入受控恢复路径。

验证范围必须覆盖真正改变执行行为的对象。只签名内核而允许替换 DTB、内核命令行、initramfs、可加载模块或策略文件，会留下等价绕过路径。镜像 manifest 应声明产品兼容范围、组件摘要、security version、签名算法和密钥标识，并由启动阶段拒绝未知或矛盾组合。

\subsection{Verified Boot 与 Measured Boot}

Verified Boot 在执行前验证签名或摘要，不满足策略则拒绝启动。
Measured Boot 把组件度量扩展到受保护寄存器或日志中，供远程或本地证明使用；它本身不一定阻止执行。
二者解决的问题不同，可以组合使用。

度量日志本身也需要完整性和事件语义。验证端必须知道每个 measurement 对应哪个阶段、算法和配置，不能只比较一个未解释的摘要。允许集合应随正式发布更新，并保留撤销已知脆弱基线的能力。

\subsection{防回滚}

合法签名的旧版本仍可能包含已知漏洞。防回滚机制应比较受保护的安全版本计数器，并只允许单调升级。
更新失败回退与安全防回滚存在张力：系统可以回退到上一个已确认且未低于安全下限的版本，但不能无限回退到任意旧镜像。

安全版本不一定等同于产品版本。功能版本可以因现场兼容性回退，安全计数器却只能在确认新版本可恢复后推进。过早烧写不可逆计数会让失败更新无法回退；过晚推进又可能留下旧漏洞窗口，因此顺序必须进入掉电状态机和目标板测试。

\subsection{恢复路径}

恢复固件和恢复密钥同样属于信任链，不能成为绕过验证的后门。
平台固件韧性应覆盖保护、检测和恢复三个能力\cite{nist-platform-resiliency}：保护关键固件免受未授权修改，检测损坏，并从独立可信副本恢复。

恢复模式应限制网络、外设和命令集合，只接受经认证的恢复对象，并清楚区分“修复系统软件”与“清除用户数据”。现场无法自动恢复时，设备至少要输出不泄露秘密的稳定故障码，使服务端能够区分签名错误、版本拒绝、介质损坏和硬件不匹配。

\section{密钥、证书与时间依赖}

密钥清单应记录用途、算法、长度、生成位置、存储边界、允许调用者、轮换方式、撤销方式、备份策略和销毁条件。只维护一张证书到期表不足以管理隐藏在 Boot ROM、云端 HSM、生产工位、恢复镜像和第三方服务中的信任关系。

根密钥宜离线或由多方控制，日常发布使用权限更小的在线或短期密钥。设备端应支持在旧密钥仍有效时安装新信任锚，并验证从旧根到新根的授权链。轮换演练必须包含旧密钥泄露、设备长期离线、证书已过期、时间不可信和部分机群升级失败等情况。

证书验证依赖时间时，要定义启动早期和长期离线设备的策略。接受任意未来时间会破坏有效期检查，时间回退又可能复活过期凭据。可采用安全保存的最近可信时间、签名时间元数据、短期会话凭据和在线状态确认组合，但每种宽限都必须有上限和审计。

密钥删除必须验证副本范围。删除文件名不表示密钥已从 FTL、备份、core dump、诊断包、CI artifact 或日志中消失。高价值私钥应尽量不可导出；普通存储中的密钥材料应使用可销毁的密钥加密密钥封装。

\section{运行时隔离}

\subsection{特权分离}

应用应以最小用户、最小 capability 和只读文件系统运行。能够拆分的网络解析、设备控制和密钥操作不应全部集中在单一 root 进程中。
systemd sandbox、seccomp、Linux Security Modules、namespace 和 cgroup 可以限制进程能力，但策略必须通过功能和攻击测试验证。

权限设计应从调用图出发：哪个进程必须打开设备节点、绑定低端口、修改网络配置或请求签名。初始化阶段临时需要的权限应在完成后主动丢弃，辅助工具不应因为“便于诊断”继承主服务全部权限。容器提供打包边界，不自动等于安全边界；共享内核、特权挂载和设备直通仍需单独评估。

编译和内核加固选项只能提高利用成本，不能修复逻辑授权错误。栈保护、FORTIFY、PIE/ASLR、只读重定位、控制流保护和内核自保护应与平台性能、调试和兼容性共同验证，并在发布 manifest 中记录实际启用状态。

\subsection{TEE 与安全服务}

Arm TrustZone 和可信执行环境可以隔离密钥、加密操作和安全存储，但不会自动保护普通世界中的协议和业务逻辑。
安全世界接口应尽量窄小，对输入长度、共享内存、命令状态和重放进行严格验证。

TEE 与普通世界之间的命令 ABI 应像网络协议一样版本化。安全服务不能信任普通世界提供的指针、长度、身份字符串和时间；返回错误也不应泄露密钥存在性、签名中间状态或可用于侧信道的细粒度差异。

\subsection{存储保护}

安全存储需要同时考虑机密性、完整性、新鲜性和可恢复性。
仅加密文件系统不能防止回放旧数据；仅做摘要也不能防止攻击者替换摘要本身。
密钥封装应绑定设备身份和启动状态，并定义主板更换、返修和数据恢复策略。

敏感状态还要定义可用性。若受保护计数器写满、RPMB 不可用、证书分区损坏或解密失败，系统应进入有界安全模式，而不是回退到明文默认值或跳过认证。存储边界和真实掉电验证见 \mbox{第~\ref{chap:storage-reliability}~章}。

\section{默认安全与攻击面治理}

secure by default 表示产品在用户不修改高级设置时即采用最小暴露和合理保护，而不是把关闭危险功能的责任交给部署者。消费物联网基线明确反对通用默认口令，并要求漏洞报告、软件更新、敏感参数保护和最小攻击面等能力\cite{etsi-en-303-645}。

每个监听端口、无线广播、USB class、串口命令、Web 页面、RPC 方法和本地 socket 都应有所有者、用途、认证方式、权限和关闭条件。开发镜像、量产测试服务和现场诊断服务必须使用不同构建或明确的生命周期策略，不能依靠隐藏 URL、未公开命令或编译后无人记得的宏开关。

默认配置应做到：

\begin{itemize}
  \item 首次使用凭据唯一或通过受控配网建立，不使用全产品共享口令；
  \item 不需要的服务、协议版本、账号和内核接口不进入镜像或默认关闭；
  \item 管理接口只在规定网络、物理动作或短期维护窗口中开放；
  \item 恢复出厂后清除旧授权，并要求重新证明所有权；
  \item 安全降级必须显式可见，不能静默关闭证书校验或访问控制。
\end{itemize}

攻击面清单应由构建产物自动校验，例如扫描开放端口、ELF capability、SUID/SGID、设备节点权限、Web route、调试符号和内核配置。基线变化需要评审；“扫描没有发现端口”不能替代对按需启动服务、无线接口和物理端口的分析。

\section{安全更新}

安全更新必须验证离线签名，使用 TLS 保护传输，并防止冻结、回滚和错误产品安装。
根密钥应尽量离线保存，在线发布密钥权限受限且可轮换。更新元数据角色分离可以降低单一在线密钥泄露的影响\cite{tuf-spec}。

TLS 保护下载链路，但不能替代离线产物签名；CDN、代理或云端账户泄露后，设备仍应拒绝未被发布策略授权的镜像。反之，只有镜像签名而没有新鲜元数据，也可能被长期冻结在已知脆弱版本。

更新机制应支持：
\begin{itemize}
  \item 根密钥和目标签名密钥轮换；
  \item 紧急撤销特定版本或密钥；
  \item 分阶段发布、暂停和回滚；
  \item 设备端安装结果与失败原因审计；
  \item 长期离线设备重新上线后的元数据更新。
\end{itemize}

更新状态机应记录下载、校验、安装、试启动、健康确认和安全计数推进的独立状态。安全修复不能绕过普通更新的掉电保护；紧急发布仍需兼容性、恢复和最小业务验证。分阶段发布应按硬件修订、存储料号、地区和风险切片观察失败率，而不是只按随机设备百分比推进。

长期离线设备可能同时面临根元数据过期、证书失效、服务器协议升级和中间版本迁移缺失。发布系统应维护可验证的升级路径或受控恢复包，并测试从仍在支持期内的最旧现场版本直接恢复，而不是只测试相邻版本升级。

安全补丁如果改变 schema、密钥格式或安全计数器，必须与回滚策略共同设计。无法回滚的数据迁移应保留向前兼容读路径或版本化副本，避免软件回退后因数据不可解释而造成另一种拒绝服务。

\section{接口与协议加固}

所有来自网络、串口、文件、总线和共享内存的数据都应被视为不可信输入。
解析器必须检查长度、整数溢出、状态转换、资源上限和超时。二进制协议应明确字节序、版本、最大消息尺寸和未知字段处理方式。

输入校验必须发生在完成资源分配和状态改变之前。长度合法不表示语义合法；解析成功后仍要检查字段组合、权限、generation、时间窗和幂等性。分片、压缩、嵌套对象和批处理会放大内存与 CPU 消耗，必须限制总解码大小、递归深度、并发会话和单位身份的请求速率。

管理接口应使用双向认证或等价的设备身份机制，并实施最小授权、重放保护、速率限制和审计。
出厂默认口令、硬编码共享密钥和未认证维护命令不应进入量产系统。

认证成功不等于授权。设备应区分读取状态、修改配置、控制执行器、安装软件、导出诊断和开放调试等权限，并把授权决定绑定租户、设备、对象和状态版本。高风险命令宜要求物理在场、双人审批或短期授权，执行结果应返回稳定命令 ID，避免网络重试造成重复动作。

协议安全还依赖错误路径。认证失败、消息过期、nonce 重复、证书未知、空间不足和时钟不可信都应有有界处理；不能在失败后回退到旧协议、明文传输或匿名模式。接口 schema、兼容性和幂等设计见 \mbox{第~\ref{chap:interface-protocol-data-contracts}~章}。

\section{构建与供应链安全}

发布系统应限制谁能够修改源码、构建配置和签名策略，并把评审、CI 结果、构建产物和签名操作关联到同一版本。
依赖下载必须验证摘要或签名；构建节点不应长期持有根签名密钥。

构建来源不仅是 Git commit。编译器、容器镜像、外部 SDK、下载镜像、补丁序列、配置片段、环境变量和生成工具都会改变产物。发布 provenance 应把输入摘要、构建者身份、执行环境和输出摘要关联起来；SLSA 规范提供了供应链 provenance 与构建完整性的分级思路\cite{slsa-spec}。

签名服务应验证候选产物已经通过规定流水线，而不是对任意上传文件提供“盖章 API”。高权限发布操作宜采用短期凭据、多方批准、不可变审计和独立撤销路径。构建节点被攻陷时，团队必须能回答哪些版本、密钥和设备群受到影响。

SBOM 用于快速定位受影响组件。最小内容至少需要组件身份、版本、供应商、依赖关系、唯一标识和生成者等信息\cite{ntia-sbom-minimum-elements}。嵌入式产品还应记录静态链接库、BusyBox applet、内核配置、固件 blob、Bootloader、容器和模型运行时，不能只列包管理器看到的顶层包。

SBOM 不等于漏洞结论。扫描命中后还要确认功能是否编译、易受攻击代码是否可达、补丁是否已回移植、产品是否具有补偿控制。VEX 可以表达某个漏洞对特定产品版本的影响判断，但必须包含依据、时间和责任人，不能用“not affected”掩盖尚未分析。

\section{安全开发与变更控制}

安全活动应进入日常开发完成定义，而不是在量产前集中渗透测试。NIST SSDF 建议准备组织、保护软件、生产安全软件并响应漏洞\cite{nist-ssdf}；项目应把这些活动映射到实际仓库、CI、发布和现场团队。

关键变更至少需要：

\begin{itemize}
  \item 更新威胁模型、接口清单和安全需求影响；
  \item 对认证、解析器、密钥、启动、更新和权限策略实施独立评审；
  \item 锁定并验证新增依赖来源、许可证、维护状态和漏洞记录；
  \item 执行 secret scanning、静态分析、单元负向测试和相关模糊测试；
  \item 记录编译加固、签名、SBOM、provenance 和安全测试证据；
  \item 明确失败回滚、现场检测和修复版本策略。
\end{itemize}

代码评审不能只检查危险函数名。身份边界、授权条件、整数与长度关系、所有权、并发状态机、错误回退和日志脱敏往往比单个 API 更关键。对 C/C++ 解析边界，宜通过小型状态机、统一长度类型和受审计库减少手写复杂度；新组件可在合适场景采用内存安全语言，但 FFI、unsafe 和底层驱动边界仍需显式验证。

测试豁免、临时证书和调试开关都应有自动到期机制。仅在工单中写“发布前删除”不可靠；CI 应直接检查量产配置中不存在测试根证书、固定 nonce、宽松认证、调试 shell 和未签名加载路径。

\section{漏洞接收、评估与 PSIRT}

每个在支持期内的产品都应提供可达的安全联系渠道、加密提交方式、确认时限和披露政策。网站和产品文档可以发布 \texttt{security.txt}，RFC 9116 规定了其机器可读格式\cite{rfc9116-securitytxt}。渠道无人值守或只接受注册用户，会促使报告者公开披露而失去修复窗口。

PSIRT 或等价责任组需要维护从接收到关闭的状态机：去重和确认、复现、受影响版本识别、严重度与现场风险评估、缓解与修复、协调披露、发布监控和事后复盘。每个状态都应有责任人和时限，并保护报告者身份与未公开细节。

CVE 标识漏洞实例，CWE 描述缺陷类型，CVSS 描述技术严重度；三者都不能代替产品影响分析。团队还应结合被利用证据、设备暴露面、权限、物理后果、机群规模、更新覆盖和恢复难度确定优先级。无法立即更新时，应提供关闭功能、网络隔离、凭据轮换或云端阻断等可验证缓解措施。

第三方公告进入后，要用 SBOM 和构建记录定位所有产品分支，包括已停止常规功能开发但仍在安全支持期的版本。补丁回移植必须记录与上游差异并重新执行相关安全与回归测试；“版本号低于修复版本”不自动代表脆弱，“版本号相同”也不证明包含同一补丁。

\section{生产与返修安全}

生产流程应把通用软件烧录与设备唯一凭据注入分离。工站需要身份认证、最小权限、操作审计和失败隔离，敏感数据不得写入普通日志。

生产系统应以设备为事务单位关联芯片 UID、板卡序列号、MAC/无线地址、证书、校准记录、烧录镜像摘要、工站、操作者和最终生命周期状态。任何中途失败都应进入隔离队列，不能简单重试后生成两个身份、复用写计数或留下半锁定设备。

私钥优先在设备内部生成，只导出证书请求或公钥。确需外部注入时，应使用受控 HSM、短期会话密钥和端到端确认，避免明文经过工站磁盘、剪贴板、日志和通用消息队列。产线网络不能直接访问根签名系统，测试账户也不应具有正式设备注册权限。

返修流程必须验证设备身份和所有权。临时开放调试口时，应使用有时限、可审计的授权，并在返修结束后恢复生产安全状态。

维修授权应绑定设备、工单、允许命令、次数或有效期，并在设备端验证签名。返修工具需要区分读取非敏感诊断、导出用户数据、替换主板、重新注入身份和解除安全锁等权限。所有维修完成后都应重新执行启动验证、调试锁、证书、更新和恢复测试。

更换主板会改变硬件身份和存储密钥。流程必须明确哪些用户数据可迁移、谁批准、旧设备身份何时撤销以及云端如何避免把新旧主板同时视为同一可信设备。不能通过复制整个数据分区绕过硬件绑定和防回滚状态。

\section{隐私、审计与取证}

安全日志需要足够支持调查，但不应成为新的敏感数据库。设计时应区分设备标识、账户信息、位置、无线邻居、音视频内容、密钥材料和普通运行指标，并定义采集目的、保留时间、访问者、导出范围和删除方式。

日志应记录安全相关状态转换，例如启动验证失败、证书轮换、调试授权、权限拒绝、更新结果和恢复出厂，但不能记录私钥、完整 token、口令、未脱敏用户载荷或可重放挑战。稳定错误码和 correlation ID 比自由文本秘密转储更适合自动判定。

审计完整性需要考虑设备被攻陷后的行为。高价值事件可以追加写入受保护计数或及时上传，但网络中断和存储满不能阻止核心安全功能。日志时间还应声明可信度；未同步实时时间可结合单调时间、启动计数和服务器接收时间重建顺序。

诊断包应采用最小采集、加密封装、访问审批和自动过期。core dump、内存快照和总线 trace 可能含有密钥与用户数据，默认不应长期保留。现场观测和证据关联见 \mbox{第~\ref{chap:observability-fleet}~章}。

\section{安全退役与介质处置}

恢复出厂、所有权转移、返修清除和安全退役是不同操作。恢复出厂通常保留设备根身份与继续服务能力；所有权转移还要撤销旧用户授权；安全退役则应阻止设备重新作为正常产品注册，并处理云端、设备和备份中的全部身份副本。

退役设备应撤销云端身份并安全清除密钥和用户数据；无法可靠擦除的存储介质应采用密码学擦除或物理销毁策略。NIST SP 800-88 把介质处置区分为 clear、purge 和 destroy，并强调方法必须匹配介质与数据敏感度\cite{nist-media-sanitization}。

对带 FTL 的 eMMC、UFS、SSD 和 SD 卡，覆盖逻辑地址不能证明所有物理页已擦除。更可靠的做法是从一开始使用每设备或每用户数据加密密钥，在退役时销毁密钥并验证备份、RPMB 元数据、日志和云端副本；高敏感介质仍可能需要厂商 sanitize 能力或物理销毁。

退役证据应记录设备身份、介质、执行方法、结果、操作者和时间。批量回收时还应防止已退役设备的证书、序列号或外壳标签被移植到伪造设备。

\section{安全支持期与现场运营}

产品发布前应公开或合同化安全支持期限、更新前提和停止支持行为。支持期不能只写年份，还要说明哪些硬件修订、地区、云服务和关键第三方依赖在范围内。供应商停止提供 BSP 或模组固件时，产品团队仍需决定替换、隔离、延长维护或提前退役。

现场资产清单应能够按产品、硬件修订、软件版本、Bootloader、关键依赖、证书状态和配置切片。没有可靠清单，就无法判断漏洞影响范围、撤销错误证书或确认补丁覆盖。设备报告版本字符串还应与签名 manifest 或证明结果核对，不能完全信任普通应用上报。

安全更新运营应监控下载、校验、安装、试启动、回滚和健康确认各阶段的成功率，并允许暂停特定切片。失败设备需要有界重试、带宽保护和人工恢复路径；不能因补丁高危就无限快速重试，造成电量耗尽或网络拥塞。

停止支持前应提前通知、提供最后安全版本、导出或迁移路径、云端降级策略和退役指引。无法继续安全联网的设备应关闭高风险远程控制或限制到本地安全模式，而不是在无补丁状态下无限期保持全部能力。

\section{安全验证}

安全验证要证明安全目标和失败路径，而不只是运行扫描器。测试计划应从威胁模型派生 misuse case，覆盖正常授权、越权、重放、资源耗尽、时间异常、掉电、版本组合和物理访问。每个用例都要有可机器判断的预期结果，避免以“设备没有崩溃”作为通过标准。

静态分析、依赖扫描和配置检查适合持续执行，但其告警需要去重、定责和基线管理。动态测试应在主机 sanitizer、仿真器和目标硬件间分层：主机环境提高模糊测试吞吐，目标板验证字节序、DMA、驱动、内存限制和真实加固配置。测试自动化与故障注入方法见 \mbox{第~\ref{chap:test-automation-hil}~章}。

模糊测试需要保存 seed、语料、目标版本、覆盖信息和最小化样本。解析器修复后应把触发样本加入回归集，同时检查相邻协议状态和资源释放；只修复单个崩溃点可能留下同类整数、状态机或拒绝服务缺陷。

安全验证至少包括：
\begin{itemize}
  \item 修改启动组件、设备树、根文件系统和版本元数据，确认验证失败；
  \item 使用错误产品、过期元数据、撤销密钥和旧版本测试更新拒绝路径；
  \item 枚举网络服务、调试接口、默认凭据和高权限进程；
  \item 对协议解析、文件格式和 ioctl 执行模糊测试；
  \item 验证密钥不会出现在镜像、日志、core dump 和普通文件系统中；
  \item 验证生产、返修、恢复和退役流程的权限边界；
  \item 注入时钟回退、随机数失败、空间耗尽、证书过期和网络分区；
  \item 对调试锁、启动引脚、外部 Flash 和反向供电路径实施物理检查；
  \item 对关键安全失败执行告警、审计、密钥撤销和恢复演练。
\end{itemize}

故障注入应关注安全检查的时间和顺序，例如在签名验证与执行之间替换数据、在安全计数推进时掉电、用电压或时钟异常跳过比较，或在世界切换和 DMA 期间改变共享内存。实验必须匹配产品攻击者能力；普通消费品与高价值无人值守设备的物理攻击预算不同。

独立渗透测试适合发现团队共同假设之外的问题，但不能代替日常开发控制。测试范围应包含硬件、移动端、云 API、生产与维修工具和更新基础设施，并在修复后执行复测。发现应映射到根因和系统性改进，而不是只关闭单个报告。

\section{数字化验收矩阵}

安全验收需要把需求、场景、证据和通过条件绑定到具体发布基线。下面给出结构示例：

\begin{table}[H]
  \centering
  \caption{产品安全生命周期验收矩阵示例}
  \begin{tabular}{p{27mm}p{40mm}p{37mm}p{30mm}}
    \hline
    场景 & 条件与注入 & 关键证据 & 通过条件 \\
    \hline
    启动与恢复 & 篡改组件、旧版本、恢复介质和掉电 & 启动日志、度量与槽位状态 & 仅执行授权且可恢复版本 \\
    身份与密钥 & 轮换、过期、撤销、时间回退和克隆 & 证书链、计数与审计记录 & 无共享秘密或旧授权复活 \\
    外部接口 & 未认证、越权、重放、畸形包和资源耗尽 & 覆盖、错误码与资源曲线 & 有界拒绝且状态不改变 \\
    供应链 & 锁定依赖、重建、SBOM 和 provenance 核验 & 输入输出摘要与签名记录 & 产物可追溯且组件完整 \\
    漏洞响应 & 模拟高危公告、回移植和分批发布 & 影响清单、时限与覆盖率 & 在目标时限内完成闭环 \\
    退役处置 & 撤权、密钥销毁、介质处置与再注册 & 设备/云状态和处置记录 & 数据不可恢复且身份不可复用 \\
    \hline
  \end{tabular}
\end{table}

验收证据应关联需求 ID、威胁模型版本、硬件修订、源码与依赖摘要、构建环境、签名密钥 ID、测试工具、原始日志、偏差批准和发布版本。更换 SoC、无线模组、Bootloader、云身份服务、文件系统或生产工位，都可能使旧结论失效。

\section{安全检查清单}

\begin{itemize}
  \item 是否存在经过评审的资产、信任边界和攻击者能力模型？
  \item 安全目标是否具有明确范围、失败行为、判定器和发布门槛？
  \item 从 Boot ROM 到关键应用的每个执行阶段是否都被认证或度量？
  \item 密钥是否具有明确的生成、存储、轮换、撤销和销毁流程？
  \item 防回滚与 A/B 恢复策略是否兼容并经过断电测试？
  \item 量产设备是否关闭或认证调试接口，并移除默认凭据？
  \item 系统服务是否遵循最小权限，所有外部输入是否有资源和长度上限？
  \item 正式产物是否具有可验证 SBOM、构建 provenance 和签名审批记录？
  \item PSIRT 是否能够接收报告、定位所有受影响版本并在目标时限内处置？
  \item 生产、返修、所有权转移和退役是否防止身份复制与权限遗留？
  \item 日志和诊断包是否既支持调查，又避免泄露密钥和用户数据？
  \item 安全支持期、停止支持策略和现场机群版本清单是否明确？
  \item 独立测试是否覆盖设备、云端、移动端、工装和更新基础设施？
\end{itemize}
